\begin{appendix}{Б}{рекомендуемое}
{Обновление корневых сертификатов}\label{ROOT}

\mbox{}

Сертификат КУЦ или корневой сертификат выступает в роли точки доверия при
проверке других сертификатов, выпускаемых в ИОК. Срок действия проверяемого
сертификата не может выходить за пределы срока действия соответствующей точки
доверия. Поэтому при длительной эксплуатации ИОК требуется периодически вводить
новые точки доверия: дополнительно к действующему корневому сертификату~$C_0$
выпускать новый корневой сертификат~$C_1$ с более поздней датой окончания
действия. Речь идет об обновлении корневого сертификата. 

При обновлении должны соблюдаться следующие правила.

\begin{enumerate}
\item
Открытый ключ сертификата $C_1$ должен генерироваться заново вместе с 
соответствующим личным ключом. 

\item
Сертификат~$C_1$ должен быть выпущен не позднее чем за $t+30$ суток до 
окончания действия $C_0$. Здесь $t$~--- максимальный срок действия в сутках
сертификатов ИОК, за исключением сертификатов КУЦ и ИУЦ.
%
Дата окончания действия~$C_1$ должна быть позже даты окончания действия~$C_0$.

\item
Сразу после выпуска $C_1$ должен быть выпущен линк-сертификат~$C_{01}$. 
В линк-сертификате должны быть указаны те же идентификационные атрибуты и 
открытый ключ КУЦ, что и в~$C_1$. Линк-сертификат должен подписываться
на личном ключе КУЦ, который соответствует открытому ключу~$C_0$.
%
Дата окончания действия $C_{01}$ должна совпадать с датой окончания 
действия~$C_1$.

Стороны ИОК могут использовать линк-сертификаты для автоматизации обновления 
своих точек доверия: если $C_0$ входит в перечень точек доверия стороны и 
$C_{01}$ признан действительным относительно $C_0$, то в перечень может быть 
добавлен сертификат~$C_1$.

\item
В линк-сертификат должны быть перенесены без изменений все расширения~$C_1$, 
кроме \code{AuthorityKeyIdentifier}. 
%
Расширение \code{AuthorityKeyIdentifier}, необязательное для $C_1$, должно 
обязательно присутствовать в линк-сертификате. В нем должно быть указано 
хэш-значение открытого ключа $C_0$ (см.~\ref{FMT.Ext.SKID}).
%
Линк-сертификат не должен содержать других расширений, кроме 
\code{AuthorityKeyIdentifier} и тех, которые содержатся в~$C_1$.

\item
Сразу после выпуска $C_1$ должен быть выпущен подчиненный сертификат 
ИУЦ. 
%
Открытый ключ этого сертификата должен генерироваться заново 
вместе с соответствующим личным ключом. 
%
Дата окончания действия сертификата ИУЦ должна совпадать с датой 
окончания действия~$C_1$.

\item
В сертификатах ИУЦ с разными точками доверия должны отличаться друг от друга
точки распространения СОС, представленные в расширении
\code{CRLDistributionPoint} (см.~\ref{FMT.Ext.CDP}).
%
Рекомендуется вести учет числа выпущенных корневых сертификатов и учитывать
порядковый номер корневого сертификата при формировании соответствующей точки 
распространения.
\end{enumerate}

\end{appendix}

